\chapter{Reuniões}

\section{Reunião Geral (RGER)}
\begin{enumerate}
  \item É convocada semanalmente pelo Team Leader, a quem cabe indicar o local, o momento de realização e a ordem de trabalhos.
  \item Em caso de necessidade, pode ser convocada por um membro da Direção.
  \item Cada Departamento deve designar um membro responsável pela apresentação do respetivo desenvolvimento, sendo essa função tipicamente assumida pelo Department Leader.
  \item Cada Departamento deverá preparar previamente a sua apresentação, através do documento disponibilizado para o efeito.
  \item A apresentação de cada Departamento deve assentar em pragmatismo, objetividade e clareza.
  \item A presença dos membros é valorizada como compromisso coletivo da equipa, devendo cada membro procurar cumpri-la.
  \item Tem como objetivos:
  \begin{enumerate}
    \item Anunciar notas gerais sobre o rumo da equipa e possíveis decisões que afetem os seus membros.
    \item Promover a transferência de conhecimento transversal a toda a equipa.
    \item Dar a conhecer informação relevante para o desenvolvimento técnico e organizacional da equipa.
    \item Dar aos Departamentos a oportunidade de apresentar o trabalho desenvolvido, promovendo o reconhecimento do esforço e dedicação dos seus membros.
    \item Reunir a equipa de forma coletiva, promovendo o convívio.
  \end{enumerate}
  \item O funcionamento deverá seguir:
  \begin{enumerate}
    \item Introdução e notas gerais da responsabilidade da Board.
    \item Apresentações sucessivas da responsabilidade de cada Departamento.
    \item Qualquer outra dinâmica considerada pertinente.
  \end{enumerate}
\end{enumerate}

\section{Reunião de Departamento (RDEP)}
\begin{enumerate}
  \item Deve ser convocada e organizada pelo Department Leader.
  \item Tem como objetivos:
  \begin{enumerate}
    \item Atribuição de tarefas.
    \item Esclarecimento sobre o estado de tarefas pendentes.
    \item Reflexão geral sobre o Departamento e o respetivo desenvolvimento técnico.
  \end{enumerate}
  \item É obrigatória a presença dos membros integrantes do Departamento.
  \item Requer planeamento prévio e registo de notas, cuja responsabilidade cabe ao Department Leader.
  \item O funcionamento deverá seguir:
  \begin{enumerate}
    \item Regime semanal, em horário preferencialmente fixo, no espaço de trabalho da equipa. Para o bom funcionamento do Departamento, deve realizar-se preferencialmente no início da semana de trabalho. Pode, em alternativa, realizar-se noutro espaço, desde que tal seja comunicado previamente pelo Department Leader.
    \item Introdução de temas gerais, tipicamente resultantes da RDIR que antecede as RDEP.
    \item Revisão sobre a semana anterior:
    \begin{enumerate}
      \item Analisar as tarefas que haviam sido atribuídas.
      \item Identificar as razões pelas quais determinadas tarefas não foram concluídas e documentar os problemas enfrentados.
      \item Redistribuir tarefas não concluídas, com eventuais alterações.
    \end{enumerate}
    \item Planeamento da semana seguinte:
    \begin{enumerate}
      \item Decidir, em conjunto com os membros, as tarefas a executar, com base no planeamento existente.
      \item Definir os Reviewers para as tarefas.
      \item Descrever detalhadamente o objetivo e os requisitos de cada tarefa a executar.
    \end{enumerate}
  \end{enumerate}
  \item Cada Departamento pode definir a metodologia de trabalho que melhor se adequa às suas necessidades, sem prejuízo do cumprimento dos objetivos mínimos definidos para a RDEP.
  \item A RDEP pode integrar dinâmicas complementares dependentes da metodologia adotada pelo Departamento, incluindo, a título exemplificativo, retrospetivas, revisões técnicas, planeamento por sprints, acompanhamento por milestones ou outras práticas associadas a Scrum, Kanban, DevOps, Waterfall, Stage-Gate ou V-Model.
  \item Podem ser convocadas RDEP complementares pelo Department Leader sempre que as necessidades do Departamento ou da equipa o justifiquem.
\end{enumerate}

\section{Reunião de Acompanhamento (RACP)}
\begin{enumerate}
  \item Funciona em regime semanal, em horário preferencialmente fixo, no espaço de trabalho da equipa. Para o bom funcionamento da equipa, deve realizar-se preferencialmente entre RDEP e RGER semanais sucessivas.
  \item É convocada pelo Department Leader. Quando necessário, qualquer membro pode propor ao Department Leader a realização de uma RACP extraordinária.
  \item Cada Subsystem deverá ter a sua RACP.
  \item Os membros presentes deverão ser os que integram o Subsystem em causa e o Department Leader.
  \item O objetivo da RACP é refletir sobre o progresso desenvolvido, verificar o cumprimento dos requisitos e permitir ao Department Leader apoiar o direcionamento do trabalho.
\end{enumerate}

\section{Reunião Inter-Subsystems (RISU)}
\begin{enumerate}
  \item A sua frequência depende da necessidade e pertinência da reunião.
  \item As reuniões podem ser propostas por um Membro, sendo convocadas por um Membro da Direção.
  \item Servem para expor e debater interdependências entre Subsystems.
  \item A obrigatoriedade de presença em RISU depende do âmbito da mesma.
  \item O final de uma fase culmina com uma RISU de presença obrigatória dos membros do Subsystem, do Department Leader e de uma representação da Board, cujo debate se cinge maioritariamente ao Subsystem em causa.
  \item Um documento escrito deve ser produzido, refletindo os temas e conclusões da RISU.
\end{enumerate}

\section{Reunião Externa (REXT)}
\begin{enumerate}
  \item A sua frequência depende da necessidade e pertinência da reunião.
  \item As reuniões devem ser convocadas por um membro da Direção.
  \item Qualquer membro pode propor a convocação de uma REXT.
  \item Estas englobam reuniões com professores, alumni, outras equipas e outros.
  \item Pode ter como objetivos:
  \begin{enumerate}
    \item Obter validação sobre o desenvolvimento realizado pela equipa.
    \item Debater qualquer tema relacionado com a Formula Student.
  \end{enumerate}
  \item Um documento escrito deve ser produzido, englobando a preparação, objetivos e resultados.
\end{enumerate}

\section{Reunião com Empresa (REMP)}
\begin{enumerate}
  \item A sua frequência depende da necessidade e pertinência da reunião, bem como da disponibilidade da empresa.
  \item As reuniões devem ser convocadas pelo Sponsors Responsible, Department Leader ou membro da Direção.
  \item Pode ter como objetivos:
  \begin{enumerate}
    \item Criar possíveis pontos de parceria, institucionais ou empresariais.
    \item Apresentar resultados e/ou desenvolvimento efetuado pela equipa.
    \item Discutir detalhadamente a metodologia a seguir para o patrocínio.
  \end{enumerate}
  \item Um documento escrito deve ser produzido, englobando a preparação, objetivos e resultados.
\end{enumerate}

\section{Reunião de Finanças (RFIN)}
\begin{enumerate}
  \item Funciona em regime mensal, em horário preferencialmente fixo. Para o bom funcionamento da equipa, deve realizar-se preferencialmente no início de cada mês.
  \item Pode ter como objetivos:
  \begin{enumerate}
    \item Promover a reflexão e debate sobre o estado atual orçamental da equipa.
    \item Informar a Board e o Financial Manager sobre decisões que afetem a equipa.
    \item Promover um ambiente de discussão de temas que influenciem o rumo financeiro da equipa.
  \end{enumerate}
  \item Um documento escrito deve ser produzido, englobando a preparação, objetivos e resultados.
  \item O Financial Manager deve apresentar um documento escrito contendo o extrato bancário mensal, bem como uma justificação detalhada para cada movimento.
  \item É obrigatória a presença de todos os membros da Board e Financial Manager.
  \item A confidencialidade dos temas debatidos deve ser assegurada.
  \item Um membro da Board ou o Financial Manager tem direito a convocar uma RFIN extraordinária sempre que considerar necessário.
\end{enumerate}

\section{Reunião de Direção (RDIR)}
\begin{enumerate}
  \item Funciona em regime semanal, em horário preferencialmente fixo. Para o bom funcionamento da equipa, deve realizar-se preferencialmente antes das RDEP semanais.
  \item Pode ter como objetivos:
  \begin{enumerate}
    \item Promover a reflexão e debate sobre o rumo da equipa.
    \item Informar a Direção sobre decisões que afetem a equipa.
    \item Promover um ambiente de discussão de temas que influenciem um ou vários Departments.
    \item Tomar decisões que afetem o desenvolvimento da equipa.
  \end{enumerate}
  \item Um documento escrito deve ser produzido, englobando a preparação, objetivos e resultados.
  \item É obrigatória a presença de todos os membros da Direção.
  \item A confidencialidade dos temas debatidos deve ser assegurada.
  \item Um membro da Direção tem direito a convocar uma RDIR extraordinária sempre que considerar necessário.
\end{enumerate}

\section{Reunião de Board (RBOR)}
\begin{enumerate}
  \item Funciona em regime semanal, em horário preferencialmente fixo. Para o bom funcionamento da equipa, deve realizar-se preferencialmente antes da RDIR semanal.
  \item Pode ter como objetivos:
  \begin{enumerate}
    \item Promover a reflexão e debate sobre o rumo da equipa.
    \item Tomar decisões que influenciem a equipa de um ponto de vista global (e.g. orçamental, competições).
    \item Preparar a exposição a realizar em RDIR.
  \end{enumerate}
  \item Um documento escrito deve ser produzido, englobando a preparação, objetivos e resultados.
  \item É obrigatória a presença de todos os membros da Board.
  \item A confidencialidade dos temas debatidos deve ser assegurada.
  \item Um membro da Board tem direito a convocar uma RBOR extraordinária sempre que considerado necessário.
\end{enumerate}
